\section{Identifying the preferred longitude}\label{me:geometric}
The preceding theorems construct an analytic characteristic from
an actual source and a specified recurrence. Its geometric
identification uses two separately proved classical statements.
For all-circle trees, \cite[Theorem 1.2 and Corollary 7.1]{Nonparabolic}
construct the full original-group filled germ at its positive
hyperbolic endpoint, with the $M^{-1}$ flag and primitive
logarithms. They give uniqueness after the actual endpoint and
those marks have been identified. For strict sources with bare
junctions, \cite[Theorem 1.3]{Longitudes} identifies the stationary
branch of precisely \eqref{me:deformed-action}, retaining the
original physical edges and all bare junctions. Its proof uses
classical zero-circle insertion and elimination; it includes the
half-label fold and makes no quantum asymptotic assumption.

We spell out the all-circle interface to avoid replacing a group
identification by an equality of second derivatives. Put
$Q_v=e^{\ii\tau_v}$ and $E_e=e^{2\pi\ii x_e}$. Differentiating
\eqref{me:local-actions} and using \eqref{me:sine-ratios}, the
wire stationary equation is exactly
\begin{equation}\label{me:edge-equation}
 (E_e+Q_v)(E_e+Q_v^{-1})(E_e+Q_w)(E_e+Q_w^{-1})
 =(E_e-1)^2(1-E_eM^2)(1-E_e/M^2).
\end{equation}
The local source log satisfies
\begin{equation}\label{me:angle-log}
 -2\pi b_\tau(x,\tau)=
          \log\frac{Q+E}{1+QE}.
\end{equation}
To check its branch as well as its derivative, differentiate both
sides in $x$: both derivatives are
$2\pi\sin\tau/(\cos2\pi x+\cos\tau)$; both sides vanish at
$x=0$. Continue from the same complete chart. The original angle
stationarity therefore gives
\begin{equation}\label{me:primitive}
 \sum_{e\ni v}\log\frac{Q_v+E_e}{1+Q_vE_e}
       -n_v\log Q_v=-2\pi\ii,
 \qquad n_v=\epsilon_v-2p_v.
\end{equation}
These are exactly $U_v+p_vV_v=-2\pi\ii$, with
$U_v=\sum\log((Q_v+E_e)/(1+Q_vE_e))-\epsilon_v\log Q_v$
and $V_v=2\log Q_v$. The linear first-charge term in
\eqref{me:base-action} fixed the right side before exponentiation.

At $t=0$ the selected source point is the positive smooth endpoint
identified by the original summed or directed source theorem;
\cite[Section 8]{Directed} includes its backwards restoration from
conditional saddles. The summed criterion gives
$\sqrt{d_v^2+n_v^2}\ge\sqrt{2^{25}d_v}>630$ at every vertex,
and the directed criterion gives
$\sqrt{d_v^2+n_v^2}=2\pi|D_v|\ge2\pi\,2^{15}>630$.
Thus the marked geometric theorem applies also to high-degree
vertices with zero twist. Its full-group reconstruction uses
\eqref{me:edge-equation} and the original flags, not just traces.
At that endpoint $Q_v^2\ne1$ for each filled disk. It persists on
a neighborhood, so a slope word commuting with its semisimple
longitude and having eigenvalue ratio one is projectively identity.
Consequently the reconstructed branch is a branch of the actual
filled group. The endpoint, flag and primitive logarithmic
conditions now meet the hypotheses of
\cite[Corollary 7.1]{Nonparabolic}, which identifies the branch.

For the geometric conclusion in the strict class, enlarge its
fixed power threshold to the maximum of the analytic threshold in
\cite[Theorem 1.2]{MovingSources} and the classical positive-filling
and longitude threshold in \cite[Theorem 1.3]{Longitudes}. Both
are constants of the same fixed original diagram, so this preserves
independent signed powers and imposes no ratio condition. The
analytic theorems above themselves require only their analytic
source threshold. No inequality between the two thresholds is
being assumed.

For a strict source beyond this common threshold the classical action in
\cite[Theorem 1.3]{Longitudes} is the sum of the same $m$ and $b$
functions, with their absolute primitives, the same chirp and the
same first-charge term. Hidden unturned wires can be retained
before its stationary elimination. Its original physical-edge
set is $\mathcal E_T$. The nondegenerate source branch proved
here therefore equals that classical branch by the local implicit
uniqueness theorem; no individual half-edge elimination is used
at a folded label.

\begin{corollary}[The selected geometric factor]\label{me:geometric-factor}
For every fixed source in Section~\ref{me:intro}, with the strict-class
threshold enlarged to the common threshold just specified, the germ
$L_{\rm an}$ is its preferred-longitude eigenvalue on the $M$
line, equivalently the inverse value on the source $M^{-1}$ line.
It has lift $L_{\rm an}(1)=-1$. Every specified actual all-level
annihilator in Theorem~\ref{me:symbol} contains this selected
geometric factor in its first nonzero classical symbol. If
$\tau_K=\left.\frac{d}{du}\log[-L_{\rm an}(e^u)]\right|_{u=0}$,
then uniformly on the stated charts,
\[
 \frac{F_N(q_N^s)}{K_NH_N(s)}
       =1+\frac{\pi\ii\tau_K s^2}{2N}+O(N^{-2}).
\]
\end{corollary}
\begin{proof}
The branch identification just established permits the marked
longitude formulas \cite[Theorems 1.2--1.3]{Longitudes}. Their
rational product is \eqref{me:longitude-formula}, including its
integer exponent $a(D)$ and negative determinant-one sign. Those
constants are fixed there by peripheral and framing calculations,
not by the quantum tangent. Hence $\tau_K=\mathcal S''(0)
=S''(0)/(\pi\ii)$. Theorem~\ref{me:symbol} proves the factor
statement, and \eqref{me:first-coefficient} gives the expansion.
\end{proof}

\section{Finite reproduction and boundaries}\label{me:checks}
The accompanying \texttt{check.py} checks the finite rational
reduced/full gauge for several operator orders, the rational sine
ratios in \eqref{me:sine-ratios}, and coefficient recovery and
finite-jet truncation for sparse formal polynomials. These are
algebraic checks of displayed identities. They do not establish
analytic remainders, global source gaps, existence of a recurrence,
or a geometric branch by testing levels. Those are the written
proofs and explicitly named companion inputs above.

The expansion is microscopic on each fixed compact in $s$.
The auxiliary $t$ family organizes its local coefficients and
creates an analytic germ, but does not supply a global asymptotic
for fixed $M\ne1$. The recurrence hypothesis is one fixed finite
operator valid for the whole function at all sufficiently large
levels. The conclusion selects one germ and one irreducible factor;
it does not assert equality with a full $A$-polynomial or an
operator ideal. None of these distinctions alters the original
all-level and both-parity source conclusions.
